\markboth{DATA SOURCES AND INTEGRATION}{DATA SOURCES AND INTEGRATION}
\section{Data Sources and Integration}\label{sec:data}

Training Copilot integrates seven data sources through the uniform MCP interface described in Section~\ref{sec:architecture:mcp}. Each is encapsulated in an independent server handling authentication, rate limiting, error recovery, and data normalisation internally; Table~\ref{tab:data_summary} summarises sources, tool counts, and metrics. Every server here is our own native FastMCP implementation, each wrapping a different kind of upstream behind clean, uniform MCP tools: Strava on its official v3 REST API, Garmin on the open-source \texttt{garminconnect} library---which reverse-engineers Garmin Connect's undocumented internal web endpoints, since Garmin exposes no public wellness API---weather on Open-Meteo, routes on OpenRouteService together with Google Geocoding and OSM Overpass, and calendar on the Google Calendar API; the flythrough render server (Section~\ref{sec:architecture:supporting}) is likewise our own code. The one deliberate exception is the Telegram bridge (Section~\ref{sec:data:telegram}), which is \textit{not} a server we wrote but an external, prefabricated MCP project vendored unchanged and merely proxied onto our transport---included precisely to demonstrate that the architecture ingests a third-party server as readily as our first-party ones. Every server runs as an independent process over Streamable HTTP and, in the containerised deployment, as its own Docker container (one per MCP server, alongside one per agent) built from a single shared Dockerfile (Section~\ref{sec:architecture:deployment}).\footnote{For all development and evaluation we used the authors' own Strava and Garmin accounts and personal training data; access to these accounts can be provided to graders on request.} This section describes the integration considerations for each; Table~\ref{tab:api_overview} in Appendix~\ref{sec:appendix:apis} complements it with a per-API reference of access requirements (API keys, OAuth, credentials), cost/rate limits, and native MCP availability for every external service used, including the LLM backend.

\subsection{Strava --- Activity Tracking}\label{sec:data:strava}

The Strava server (port 8103) connects to the Strava v3 REST API via the standard OAuth\,2.0 Authorization Code flow: on first use, the server opens Strava's consent page, the user grants the requested scopes (\texttt{read}, \texttt{activity:read\_all}, \texttt{activity:write}), and the resulting authorization code is exchanged for access and refresh tokens, persisted to \texttt{.tokens/strava.json} and refreshed automatically with a five-minute buffer before reported expiry. This flow (\texttt{auth/strava\_oauth.py}) is transparent to the MCP tools, which simply call the API with a valid token.

Its thirteen tools cover recent activities, aggregate statistics, year-over-year breakdowns, personal records, gear mileage, deep single-activity detail with lap splits and power data, GPS streams at sub-second resolution, performance-trend analysis against personal baselines, and the Training Stress Balance metrics (CTL, ATL, TSB) derived from the fitness--fatigue impulse-response model~\citep{morton_modeling_1990}. HTTP errors use exponential backoff for server errors (5xx) but return immediately on rate limits (429), since Strava's \texttt{Retry-After} is typically 15\,+ minutes, making in-request retries pointless. In June 2026, Strava announced that API access now requires a subscription for Standard Tier developers, alongside an official Strava MCP server for AI-native access~\citep{strava_api_update_2026}---validating Training Copilot's MCP-first architecture, since the official server could eventually replace ours with a registry URL change and no agent or UI code changes. Figure~\ref{fig:screenshot_activity} in Appendix~\ref{sec:appendix:screenshots} shows the resulting Dashboard tab, which combines an activity map, recent-activity list, and training-load trend charts drawn from these tools. A dedicated \textbf{Activity Analysis} view renders per-stream charts (heart rate, pace, elevation, cadence, power) alongside an interactive route map that recolours by the selected metric, turning a single activity's raw GPS and sensor streams into a diagnostic overlay. A companion \textbf{Sync tab} (Figure~\ref{fig:screenshot_sync}) transfers activities from Garmin to Strava without the user manually downloading and re-uploading files: it fetches Garmin activities for a chosen date range and shows a \textit{preview}---each activity as a card with its route map and a flag for whether it is already on Strava---from which the user selects which to transfer before the tool downloads each Garmin FIT file and uploads it directly to Strava. This helps athletes who record on a Garmin watch but keep their social and analysis history on Strava.

\subsection{Garmin Connect --- Wellness and Recovery}\label{sec:data:garmin}

Garmin provides no public REST API for the wellness data (sleep, HRV, Body Battery, stress) Training Copilot needs---the Garmin Health API is enterprise-only---so our server relies on the open-source \texttt{garminconnect} Python library,\footnote{\url{https://github.com/cyberjunky/python-garminconnect}} which reverse-engineers the Garmin Connect web app's internal endpoints. Authentication uses Garmin's SSO login with stored credentials and MFA on first login; session tokens are cached in \texttt{.tokens/} and reused until expiry. This carries inherent risk: the endpoints are undocumented and may change without notice, and the library has no official support. Rate limits are stricter and less predictable than Strava's---the Body Battery endpoint rejects date ranges over 28 days, so the server chunks requests automatically---and a retry mechanism with exponential backoff handles transient failures. A mock-data mode (\texttt{GARMIN\_MOCK\_HEALTH=true}) enables UI development without a connected device.

Its thirteen tools span seven categories: sleep stages with SpO\textsubscript{2} and sleep-associated HRV; last-night HRV with personal baseline and readiness status; Body Battery as daily highs, lows, and intraday timeline; intraday stress and heart-rate timelines; VO\textsubscript{2}max, training load, and race predictions; body composition; and multi-day wellness trends with step timelines and activity GPS tracks. These metrics let the recovery specialist assess training readiness by jointly considering sleep quality, HRV relative to baseline, stress, and remaining energy---the multi-signal approach \cite{rothschild_predicting_2024} showed is necessary for accurate recovery prediction. Figure~\ref{fig:screenshot_health} in Appendix~\ref{sec:appendix:screenshots} shows the resulting Health tab.

\subsection{Weather, Routes, and Calendar}\label{sec:data:other}

The \textbf{weather server} (port 8101) queries the Open-Meteo API, which needs no API key and imposes no rate limits for non-commercial use. Its four tools provide current conditions, multi-day forecasts (1--16 days, hourly or daily) with temperature, precipitation probability, wind speed, and UV index, pollen concentrations by species, and UV category classification; the context specialist applies heuristic rules to classify training windows---5--20\,\textdegree C as ideal, precipitation probability above 30\,\% as a rain warning, UV index above 6 as a sun-protection advisory.

The \textbf{routes server} (port 8102) aggregates three external services behind seven tools. OpenRouteService (via \texttt{ORS\_API\_KEY}) provides A-to-B routing across profiles (cycling, running, hiking), circular route generation from a start point and target distance, elevation profiles, and isochrone maps; Google Geocoding resolves place names to coordinates; OSM Overpass retrieves boundary polygons for named areas, enabling park-constrained loops (a ``run inside Schlossgarten'' query generates the park boundary from OpenStreetMap, then routes within it). Each tool returns GeoJSON-compatible coordinate arrays rendered as interactive Folium maps; Figure~\ref{fig:screenshot_routes} in Appendix~\ref{sec:appendix:screenshots} shows the resulting Routes tab.

The \textbf{calendar server} (port 8105) provides six tools for Google Calendar read/write access, supporting both a single-user development mode (persisted OAuth tokens in \texttt{.tokens/google.json}) and a multi-tenant mode (per-request \texttt{Authorization: Bearer} headers). The context specialist cross-references free calendar windows with weather forecasts to suggest optimal training times and can create events for planned sessions.

\subsection{Fitness Literature --- RAG Knowledge Base}\label{sec:data:rag}

The fitness specialist (port 9005) is the only agent without an MCP server: it queries a local vector database of fifty public-domain books on physical culture, athletics, physiology and health from Project Gutenberg\footnote{\url{https://www.gutenberg.org/}} via RAG~\citep{lewis_retrieval-augmented_2020}. Embeddings come from the local \texttt{all-MiniLM-L6-v2} model ($\sim$90\,MB)~\citep{reimers_sentence-bert_2019}; the corpus is split offline into $\sim$900-character chunks with 150-character overlap ($\sim$32k chunks in total), and each query retrieves its top five chunks by brute-force cosine similarity---a single NumPy matrix--vector product that needs no FAISS index or vector-database service at this corpus size. The specialist is instructed to ground every claim in retrieved passages with explicit source citations, and an \texttt{ensure\_sources()} post-processing step guarantees cited references survive the orchestrator's synthesis. Figure~\ref{fig:rag_pipeline} illustrates the pipeline.

\begin{figure}[ht]
\centering
\resizebox{\textwidth}{!}{%
\begin{tikzpicture}[
    node distance=0.4cm and 0.5cm,
    every node/.style={font=\small},
    block/.style={draw, thick, rounded corners=3pt, minimum height=0.7cm, minimum width=2.0cm, align=center},
    data/.style={draw, thick, rounded corners=1pt, minimum height=0.5cm, align=center, font=\scriptsize},
    arr/.style={-{Stealth[length=5pt]}, thick},
]
    \node[block, draw=tcblue, fill=tcblue!8] (query) {User Query};
    \node[block, draw=tcgreen, fill=tcgreen!8, right=1.0cm of query] (embed) {Embed\\{\scriptsize all-MiniLM-L6-v2}};
    \node[block, draw=tccyan, fill=tccyan!8, right=1.0cm of embed] (store) {Cosine Search\\{\scriptsize (N$\times$384 float32)}};
    \node[block, draw=tcyellow!80!black, fill=tcyellow!8, right=1.0cm of store] (topk) {Top-$k$\\Passages};
    \node[block, draw=tcblue, fill=tcblue!8, below=1.0cm of topk] (llm) {ReAct Agent\\{\scriptsize (LLM synthesis)}};
    \node[block, draw=tcgreen, fill=tcgreen!8, below=1.0cm of llm] (answer) {Grounded Answer\\+ Sources};
    \node[data, draw=tcgrey!40, fill=tcgrey!4, below=0.4cm of store] (corpus) {8 books, $\sim$3k chunks};

    \draw[arr, color=tcgrey] (query) -- (embed);
    \draw[arr, color=tcgreen] (embed) -- node[above, font=\scriptsize]{384-d vector} (store);
    \draw[arr, color=tccyan] (store) -- node[above, font=\scriptsize]{top 5} (topk);
    \draw[arr, color=tcyellow!80!black] (topk) -- (llm);
    \draw[arr, color=tcgreen] (llm) -- node[right, font=\scriptsize]{+ ensure\_sources()} (answer);
    \draw[arr, dashed, color=tcgrey!40] (corpus) -- (store);

\end{tikzpicture}%
}% end resizebox
\caption[RAG pipeline of the fitness specialist]{RAG pipeline of the fitness specialist. The user query is embedded locally, compared against the normalised chunk vectors via dot product, and the top-$k$ passages are fed to the ReAct agent for synthesis. The \texttt{ensure\_sources()} step guarantees that book citations survive the orchestrator's synthesis.}\label{fig:rag_pipeline}
\end{figure}


\subsection{Telegram --- External Server Bridge}\label{sec:data:telegram}

An optional Telegram integration (port 8106) demonstrates the architecture's ability to bridge external stdio-only MCP servers onto the Streamable HTTP bus without forking their code. The upstream \texttt{chigwell/telegram-mcp} project\footnote{\url{https://github.com/chigwell/telegram-mcp}} (116 tools) runs unmodified via \texttt{uv} in an isolated environment; to \texttt{ToolHost} it looks identical to every other server, confirming the vendor-agnostic promise.

\subsection{Summary}\label{sec:data:summary}

\begin{table}[ht]
\begin{center}
\begin{minipage}{\textwidth}
\caption[Summary of integrated data sources]{Summary of integrated data sources.}\label{tab:data_summary}
\begin{tabularx}{\textwidth}{llcX}
\toprule
\textbf{Source} & \textbf{Server} & \textbf{Tools} & \textbf{Primary Metrics} \\
\midrule
Strava & :8103 & 13 & Activities, load, trends, splits, GPS, records, analysis \\
Garmin & :8104 & 13 & Sleep, HRV, Body Battery, stress, VO\textsubscript{2}max \\
Weather & :8101 & 4 & Forecast, UV, pollen, current conditions \\
Routes & :8102 & 7 & A-B routes, loops, trails, isochrones, elevation \\
Calendar & :8105 & 6 & Events, free slots, event detail, create/update/delete \\
Literature & RAG & 1 & Exercise-science passages (8 books) \\
Telegram & :8106 & 116 & Messages, chats, contacts, media \\
\botrule
\end{tabularx}
\end{minipage}
\end{center}
\end{table}


Table~\ref{tab:full_tools} in Appendix~\ref{sec:appendix:tools} expands this summary into the complete tool-by-tool inventory across all six native MCP servers (the external Telegram bridge's 116 tools are summarised in Table~\ref{tab:data_summary} rather than enumerated individually).
